\begin{appendices}
\section{Two-Tower Model Implementation}
\label{app:two_tower_model}
\addcontentsline{toc}{section}{Two-Tower Model Implementation}

The complete PyTorch implementation of the two-tower neural network model for job recommendation is presented below. This code implements both the user and job towers with their embedding layers, fully-connected layers, batch normalization, and regularization techniques described in Chapter \ref{chap:user_recommendation}.

\begin{listing}[h]
\begin{minted}[fontsize=\footnotesize]{python3}
import torch
import torch.nn as nn
import torch.nn.functional as F

class TwoTowerModel(nn.Module):
    def __init__(self, num_user_cat, num_job_cat, num_user_numeric, num_job_numeric, output_dim=64):
        super().__init__()
        # User tower embeddings and layers
        self.user_embeddings = nn.ModuleList([nn.Embedding(card, 16) for card in num_user_cat])
        user_input_dim = len(num_user_cat) * 16 + num_user_numeric
        self.user_tower = nn.Sequential(
            nn.Linear(user_input_dim, 512),
            nn.BatchNorm1d(512),
            nn.ReLU(),
            nn.Dropout(0.3),
            nn.Linear(512, 256),
            nn.BatchNorm1d(256),
            nn.ReLU(),
            nn.Dropout(0.2),
            nn.Linear(256, output_dim)
        )

        # Job tower embeddings and layers
        self.job_embeddings = nn.ModuleList([nn.Embedding(card, 16) for card in num_job_cat])
        job_input_dim = len(num_job_cat) * 16 + num_job_numeric 
        self.job_tower = nn.Sequential(
            nn.Linear(job_input_dim, 512), 
            nn.BatchNorm1d(512),
            nn.ReLU(), 
            nn.Dropout(0.3),
            nn.Linear(512, 256),
            nn.BatchNorm1d(256),
            nn.ReLU(),
            nn.Dropout(0.2),
            nn.Linear(256, output_dim)
        )

    def forward(self, u_cat, u_num, j_cat, j_num):
        # Process user tower
        if len(self.user_embeddings) > 0:
            u_embeds = torch.cat([emb(u_cat[:, i]) for i, emb in enumerate(self.user_embeddings)], dim=1)
            user_input = torch.cat([u_embeds, u_num], dim=1)
        else:
            user_input = u_num
        user_vec = self.user_tower(user_input)

        # Process job tower
        if len(self.job_embeddings) > 0:
            j_embeds = torch.cat([emb(j_cat[:, i]) for i, emb in enumerate(self.job_embeddings)], dim=1)
            job_input = torch.cat([j_embeds, j_num], dim=1)
        else:
            job_input = j_num
        job_vec = self.job_tower(job_input)

        # Return matching score as dot product
        return (user_vec * job_vec).sum(dim=1)
\end{minted}
\caption{Complete Two-Tower Model Implementation in PyTorch}
\label{lst:two_tower_complete}
\end{listing}

The model takes four inputs:
\begin{itemize}
    \item \textbf{u\_cat}: User categorical features (shape: batch\_size $\times$ num\_user\_cat)
    \item \textbf{u\_num}: User numeric features (shape: batch\_size $\times$ num\_user\_numeric)
    \item \textbf{j\_cat}: Job categorical features (shape: batch\_size $\times$ num\_job\_cat)
    \item \textbf{j\_num}: Job numeric features (shape: batch\_size $\times$ num\_job\_numeric)
\end{itemize}

The forward pass produces a matching score for each user-job pair, which is used for ranking and binary classification during training.

\newpage
\section{Additional Dataset Statistics}
\label{app:dataset_statistics}
\addcontentsline{toc}{section}{Additional Dataset Statistics}
\end{appendices}

%%%%%%%%%%%%%%%%%%%%%%%%%%%%%%
%    Dutch version           %
%%%%%%%%%%%%%%%%%%%%%%%%%%%%%%
% \begin{appendices}
% \section*{Bijlage A}
% \addcontentsline{toc}{section}{Bijlage A}  

% Toelichting bijlage.



% \newpage
% \section*{Bijlage B}
% \addcontentsline{toc}{section}{Bijlage B}  

% Toelichting bijlage.

% \end{appendices}
